% \iffalse meta-comment
%
% registerstyle.dtx
% Copyright 2026 James P. Howard, II <jh@jameshoward.us>
%
% This work may be distributed and/or modified under the
% conditions of the LaTeX Project Public License, either version 1.3c
% of this license or (at your option) any later version.
% The latest version of this license is in
%   https://www.latex-project.org/lppl.txt
% and version 1.3c or later is part of all distributions of LaTeX
% version 2008-05-04 or later.
%
% This work has the LPPL maintenance status `maintained'.
% The Current Maintainer of this work is James P. Howard, II.
%
% This work consists of the file registerstyle.dtx
% and the derived file          registerstyle.sty.
%
%<*driver>
\documentclass{ltxdoc}
\usepackage[T1]{fontenc}
\usepackage{lmodern}
\usepackage{hyperref}
\usepackage{enumitem}
\EnableCrossrefs
\CodelineIndex
\RecordChanges
\begin{document}
\DocInput{registerstyle.dtx}
\end{document}
%</driver>
% \fi
%
% \changes{v1.0}{2026/07/27}{First public release.}
%
% \GetFileInfo{registerstyle.sty}
%
% \title{The \textsf{registerstyle} package\thanks{This document
%   corresponds to \textsf{registerstyle}~v1.0,
%   dated~2026/07/27.}}
% \author{James P. Howard, II\\ \texttt{jh@jameshoward.us}}
% \date{2026/07/27}
% \maketitle
%
% \begin{abstract}
% The \textsf{registerstyle} package provides commands and environments
% for typesetting genealogical registers in the style used by the
% \emph{New England Historical and Genealogical Register} (NEHGS).
% It manages a person registry with computed generation numbers,
% ancestor chains, child lists with continuation markers, multiple
% marriages, vital records, date formatting, gender-aware pronouns,
% source citations, index generation, and cross-references.
% \end{abstract}
%
% \tableofcontents
%
% \section{Introduction}
%
% The Register style, also called the NEHGS style, is the standard
% format for presenting genealogical data in descending order.
% Each person who continues the lineage receives a numbered main entry;
% children are listed under their parents with roman numeral labels.
% Children who appear later as main entries are marked with a ``+''
% continuation symbol.
%
% This package automates the mechanical parts of Register typesetting:
% numbering, generation computation, ancestor chains, parent
% validation, and index entries.  The genealogical prose itself
% remains the author's responsibility.
%
% \section{Quick start}
%
% \begin{verbatim}
% \documentclass{article}
% \usepackage[en-US]{datetime2}
% \usepackage[date=en-US,abbrevparents]{registerstyle}
%
% \definepersonkv{john-doe}{
%   name        = {John Doe},
%   preferred   = {John},
%   parents     = {},
%   sex         = {M},
%   birth       = {1820-03-15},
%   birth-place = {Boston, Massachusetts}
% }
%
% \definepersonkv{mary-smith}{
%   name        = {Mary Smith},
%   preferred   = {Mary},
%   parents     = {},
%   sex         = {F},
%   birth       = {1825-06-20}
% }
%
% \definepersonkv{james-doe}{
%   name        = {James Doe},
%   preferred   = {James},
%   parents     = {john-doe,mary-smith},
%   lineage-parent = {john-doe},
%   sex         = {M},
%   birth       = {1850-11-02}
% }
%
% \addmarriage{john-doe}{
%   spouse = mary-smith,
%   date   = 1845-09-10,
%   place  = {Springfield, Massachusetts}
% }
%
% \begin{document}
% \begin{registersection}
%
% \generationheading{john-doe}
% \mainperson{john-doe} was born
% \getpersonbirthdate{john-doe} in
% \getpersonbirthplace{john-doe}.
% \HeSheThey{john-doe} married
% \getpersonname{mary-smith} on
% \getmarriagedate{john-doe}{1} in
% \getmarriageplace{john-doe}{1}.
%
% Children of John Doe and Mary Smith:
% \begin{childrenof}{john-doe}{mary-smith}
%   \childref{james-doe}{b.~\getpersonbirthdate{james-doe}.}
% \end{childrenof}
%
% \generationheading{james-doe}
% \mainperson{james-doe} was born
% \getpersonbirthdate{james-doe}.
%
% \end{registersection}
% \checkregister
% \end{document}
% \end{verbatim}
%
% \section{Package options}
%
% The package accepts the following options:
%
% \begin{description}[style=nextline]
% \item[\texttt{index} (default)] Load \textsf{imakeidx} and
%   generate index entries whenever a person's name is printed
%   by |\getpersonname|.
% \item[\texttt{noindex}] Suppress all index generation.
% \item[\texttt{fullparents} (default)] Print full names in
%   parenthetical ancestor chains.
% \item[\texttt{abbrevparents}] Print preferred (short) names in
%   ancestor chains.
% \item[\texttt{ngsq}] Use NGSQ (Modified Register) numbering.
% \item[\texttt{datestyle=}\meta{style}] Set the \textsf{datetime2}
%   date style at |\begin{document}|.  The synonym \texttt{date=}
%   is also accepted.
% \item[\texttt{numbering=ngsq}] Equivalent to the \texttt{ngsq}
%   option.
% \end{description}
%
% \section{Defining persons}
%
% \DescribeMacro{\definepersonkv}
% |\definepersonkv|\marg{id}\marg{key=value pairs}
%
% The primary interface for declaring a person in the registry.
% The \meta{id} is a short identifier (e.g., |john-doe|) used to
% refer to this person throughout the document.  Available keys:
%
% \begin{description}[style=nextline, font=\ttfamily]
% \item[name] Full name as it should appear in the register.
% \item[preferred] Short or preferred name for abbreviated contexts.
% \item[parents] Comma-separated list of parent ids.  Leave empty for
%   the founding generation.
% \item[index] Index entry text.  Defaults to the \texttt{name} value.
% \item[lineage-parent] Single parent id for lineage chain display.
%   When set, ancestor chains follow only this parent rather than
%   listing all parents.
% \item[sex] |M|, |F|, or empty.  Used by pronoun helpers.
% \item[birth] Birth date in ISO format (|YYYY-MM-DD|).
% \item[birth-place] Birth place as free text.
% \item[death] Death date in ISO format.
% \item[death-place] Death place as free text.
% \item[baptism] Baptism or christening date in ISO format.
% \item[baptism-place] Baptism or christening place.
% \item[burial] Burial date in ISO format.
% \item[burial-place] Burial place.
% \item[occupation] Occupation as free text.
% \item[marriage-date] Date of (first) marriage in ISO format.
%   For multiple marriages, use |\addmarriage| instead.
% \item[marriage-place] Place of (first) marriage.
% \end{description}
%
% \DescribeMacro{\defineperson}
% |\defineperson|\marg{id}\marg{parents}\marg{name}\marg{preferred}%
% \oarg{index}
%
% Legacy positional interface.  Prefer |\definepersonkv| for new work.
%
% \section{Marriage support}
%
% \DescribeMacro{\addmarriage}
% |\addmarriage|\marg{person-id}\marg{key=value pairs}
%
% Add a marriage record to a previously defined person.  Supports
% multiple marriages per person.  Available keys:
%
% \begin{description}[style=nextline, font=\ttfamily]
% \item[date] Marriage date in ISO format.
% \item[place] Marriage place.
% \item[spouse] Person id of the spouse (must be separately defined).
% \end{description}
%
% Marriages are numbered sequentially starting from~1 in the order
% they are added.  If \texttt{marriage-date} or \texttt{marriage-place}
% is given inline in |\definepersonkv|, that counts as marriage~1.
%
% \DescribeMacro{\getmarriagedate}
% |\getmarriagedate|\marg{person-id}\marg{number}
%
% Print the date of marriage \meta{number}, formatted through
% \textsf{datetime2}.
%
% \DescribeMacro{\getmarriageplace}
% |\getmarriageplace|\marg{person-id}\marg{number}
%
% Print the place of marriage \meta{number}.
%
% \DescribeMacro{\getmarriagespouse}
% |\getmarriagespouse|\marg{person-id}\marg{number}
%
% Print the name of the spouse from marriage \meta{number}
% (via |\getpersonname|, including index entry).
%
% \DescribeMacro{\getmarriagecount}
% |\getmarriagecount|\marg{person-id}
%
% Expand to the number of recorded marriages for the person.
%
% \section{Name retrieval}
%
% \DescribeMacro{\getpersonname}
% |\getpersonname|\marg{id} prints the full name in small caps and
% generates an index entry.
%
% \DescribeMacro{\getpersonnameNI}
% |\getpersonnameNI|\marg{id} prints the full name in small caps
% without an index entry.
%
% \DescribeMacro{\getpersonpreferred}
% |\getpersonpreferred|\marg{id} prints the preferred (short) name
% in small caps.
%
% \DescribeMacro{\getpersonabbrevNI}
% |\getpersonabbrevNI|\marg{id} prints the preferred name if set,
% otherwise the full name, without an index entry.
%
% \DescribeMacro{\shortname}
% |\shortname|\marg{id} is an alias for |\getpersonpreferred|.
%
% \section{Vital record accessors}
%
% Each accessor prints the stored value, or a bracketed placeholder
% with a package warning if the value is missing.
%
% \DescribeMacro{\getpersonbirthdate}
% \DescribeMacro{\getpersonbirthplace}
% \DescribeMacro{\getpersondeathdate}
% \DescribeMacro{\getpersondeathplace}
% \DescribeMacro{\getpersonbaptismdate}
% \DescribeMacro{\getpersonbaptismplace}
% \DescribeMacro{\getpersonburialdate}
% \DescribeMacro{\getpersonburialplace}
% \DescribeMacro{\getpersonoccupation}
% \DescribeMacro{\getpersonsex}
%
% |\getpersonbirthdate|\marg{id},
% |\getpersonbirthplace|\marg{id}, etc.
%
% Date accessors format the stored ISO date through |\registerdate|.
% Place and occupation accessors print the stored text verbatim.
% |\getpersonsex|\marg{id} expands to the raw stored value
% (|M|, |F|, or empty).
%
% \section{Generation computation}
%
% \DescribeMacro{\getgeneration}
% |\getgeneration|\marg{id} computes and prints the generation number.
% Persons with no parents are generation~1; others are one more than
% their highest-generation parent.  Results are cached after the first
% computation.  Cycles are detected and produce a warning.
%
% \DescribeMacro{\generationheading}
% |\generationheading|\marg{id} prints a bold heading such as
% ``Third Generation'' based on the person's computed generation.
% Supports generations 1--20 by name; higher generations use the
% numeric form.
%
% \section{Register entries}
%
% \DescribeMacro{\mainperson}
% |\mainperson|\marg{id} typesets a numbered main register entry:
% the entry number in bold, the person's name in small caps with
% a generation superscript, and the parenthetical ancestor chain.
% Each call increments the |registerperson| counter and creates
% a label |person:|\meta{id} for cross-references.
%
% \DescribeMacro{\xrefperson}
% |\xrefperson|\marg{id} prints a hyperlinked cross-reference to
% a person's main entry, showing their name and entry number.
%
% \DescribeMacro{\childname}
% |\childname|\marg{id} prints a person's name with generation
% superscript, without an index entry.  Used in child lists.
%
% \DescribeMacro{\formatparents}
% |\formatparents|\marg{id} prints the parenthetical ancestor chain.
% Called automatically by |\mainperson| but available for standalone
% use.
%
% \section{Child lists}
%
% \DescribeEnv{children}
% |children| is an enumerated list with roman numeral labels,
% accepting optional \textsf{enumitem} keys.
%
% \DescribeEnv{childrenof}
% |childrenof|\marg{parent-a}\marg{parent-b}\oarg{enumitem keys}
% is like |children| but validates that each |\child| entry has the
% expected parent pair.
%
% \DescribeMacro{\child}
% |\child|\marg{id}\marg{description} typesets a child entry.
% The description text follows a comma after the child's name.
%
% \DescribeMacro{\childref}
% |\childref|\marg{id}\marg{description} typesets a child entry
% with a ``\textbf{+}'' continuation marker, indicating that this
% child will appear later as a |\mainperson| entry.  This is a
% core convention of the Register style.
%
% \section{Approximate dates}
%
% Genealogists frequently work with uncertain dates.  These commands
% format common qualifiers used in genealogical writing:
%
% \DescribeMacro{\circa}
% |\circa|\marg{date} prints ``c.\,\meta{date}''.
%
% \DescribeMacro{\before}
% |\before|\marg{date} prints ``bef.\ \meta{date}''.
%
% \DescribeMacro{\after}
% |\after|\marg{date} prints ``aft.\ \meta{date}''.
%
% \DescribeMacro{\between}
% |\between|\marg{date1}\marg{date2} prints
% ``bet.\ \meta{date1} and \meta{date2}''.
%
% \DescribeMacro{\aboutdate}
% |\aboutdate|\marg{date} prints ``abt.\ \meta{date}''.
%
% \DescribeMacro{\probably}
% |\probably|\marg{date} prints ``prob.\ \meta{date}''.
%
% \DescribeMacro{\say}
% |\say|\marg{date} prints ``say \meta{date}''.
%
% These may be used both in prose and inside key values.
%
% \section{Gender-aware pronouns}
%
% These commands select pronouns based on the |sex| field
% (|M|\,$\to$\,masculine, |F|\,$\to$\,feminine,
% other/empty\,$\to$\,neutral).
%
% \DescribeMacro{\heshethey}
% \DescribeMacro{\himherthem}
% \DescribeMacro{\hishertheirargs}
% \DescribeMacro{\HeSheThey}
% \DescribeMacro{\HimHerThem}
% \DescribeMacro{\HisHerTheir}
% \DescribeMacro{\himselfherselfthemself}
%
% \noindent
% All take a single \marg{person-id} argument.
% Lowercase and capitalized variants are provided.
%
% \section{Source citations}
%
% \DescribeMacro{\registersource}
% |\registersource|\marg{text} produces a footnote citation.
%
% \DescribeMacro{\registerendnote}
% |\registerendnote|\marg{text} produces an endnote if the
% \textsf{endnotes} package is loaded, otherwise falls back to
% a footnote.
%
% \section{Date formatting}
%
% \DescribeMacro{\registerdate}
% |\registerdate|\marg{YYYY-MM-DD} formats an ISO-8601 date using
% the current \textsf{datetime2} style.  Non-ISO input is passed
% through with a warning.
%
% \DescribeMacro{\registerdatestyle}
% |\registerdatestyle|\marg{style} changes the \textsf{datetime2}
% date style mid-document.
%
% \section{Shorthands}
%
% \DescribeMacro{\born}
% \DescribeMacro{\died}
% \DescribeMacro{\married}
% \DescribeMacro{\baptized}
% \DescribeMacro{\buried}
% \DescribeMacro{\place}
% \DescribeMacro{\dated}
%
% \noindent
% |\born|\marg{text} prints ``b.\ \meta{text}'';
% |\died| prints ``d.'';
% |\married| prints ``m.'';
% |\baptized| prints ``bp.'';
% |\buried| prints ``bur.''.
% |\place|\marg{text} prints \meta{text} in italics;
% |\dated|\marg{text} prints \meta{text} as-is.
%
% \section{Register section}
%
% \DescribeEnv{registersection}
% The |registersection| environment sets up typographic parameters
% suitable for register prose: paragraph indentation of 1.5em,
% no paragraph skip, and compact lists.  It also defines a local
% |\generation|\marg{label} command for manual generation headings.
%
% \section{Validation}
%
% \DescribeMacro{\checkperson}
% |\checkperson|\marg{id} runs validation checks on a single person:
% verifies the name is non-empty, all listed parents are defined,
% all dates are in ISO format, and places are not given without
% corresponding dates.  Problems produce package warnings.
%
% \DescribeMacro{\checkregister}
% |\checkregister| runs |\checkperson| on every registered person.
% Call this near |\end{document}| to catch data integrity issues.
%
% \StopEventually{%
%   \PrintChanges
%   \PrintIndex
% }
%
% \section{Implementation}
%
% \subsection{Package identification and options}
%
%    \begin{macrocode}
%<*package>
%% registerstyle.sty --- NEHGS-style genealogical register
%% Copyright 2026 James P. Howard, II <jh@jameshoward.us>
%%
%% This work may be distributed and/or modified under the
%% conditions of the LaTeX Project Public License, either version 1.3c
%% of this license or (at your option) any later version.
%% The latest version of this license is in
%%   https://www.latex-project.org/lppl.txt
%% and version 1.3c or later is part of all distributions of LaTeX
%% version 2008-05-04 or later.
%%
%% This work has the LPPL maintenance status `maintained'.
%% The Current Maintainer of this work is James P. Howard, II.

\NeedsTeXFormat{LaTeX2e}[2020/10/01]
\ProvidesPackage{registerstyle}[2026/07/27 v1.0 NEHGS-style genealogical register]
%    \end{macrocode}
%
% Option booleans and key=value option parser.
%
%    \begin{macrocode}
\newif\ifr@index
\r@indextrue
\newif\ifr@option@abbrevparents
\r@option@abbrevparentsfalse
\newif\ifr@option@ngsq
\r@option@ngsqfalse
\newcommand{\r@datestyle}{}
\newcommand{\r@optionkey}{}
\newcommand{\r@optionvalue}{}
\newcommand{\r@datestylekey}{datestyle}
\newcommand{\r@datekey}{date}
\newcommand{\r@numberingkey}{numbering}
\newcommand{\r@splitoption}{}
\def\r@splitoption#1=#2\@nil{%
  \renewcommand{\r@optionkey}{#1}%
  \renewcommand{\r@optionvalue}{#2}%
}
\newcommand{\r@ngsqvalue}{ngsq}
\newcommand{\r@handleunknownoption}{%
  \@expandtwoargs\in@{=}{\CurrentOption}%
  \ifin@
    \expandafter\r@splitoption\CurrentOption\@nil
    \ifx\r@optionkey\r@datestylekey
      \renewcommand{\r@datestyle}{\r@optionvalue}%
    \else
      \ifx\r@optionkey\r@datekey
        \renewcommand{\r@datestyle}{\r@optionvalue}%
      \else
        \ifx\r@optionkey\r@numberingkey
          \ifx\r@optionvalue\r@ngsqvalue
            \r@option@ngsqtrue
          \else
            \PackageWarning{registerstyle}{Unknown numbering style `\r@optionvalue'}%
          \fi
        \else
          \PackageWarning{registerstyle}{Unknown option `\CurrentOption'}%
        \fi
      \fi
    \fi
  \else
    \PackageWarning{registerstyle}{Unknown option `\CurrentOption'}%
  \fi
}

\DeclareOption{index}{\r@indextrue}
\DeclareOption{noindex}{\r@indexfalse}
\DeclareOption{fullparents}{\r@option@abbrevparentsfalse}
\DeclareOption{abbrevparents}{\r@option@abbrevparentstrue}
\DeclareOption{ngsq}{\r@option@ngsqtrue}
\DeclareOption*{\r@handleunknownoption}
\ProcessOptions\relax
%    \end{macrocode}
%
% \subsection{Required packages}
%
%    \begin{macrocode}
\RequirePackage{xparse}
\RequirePackage{enumitem}
\RequirePackage{etoolbox}
\RequirePackage{keyval}
\RequirePackage{xspace}
\RequirePackage{datetime2}
\RequirePackage{hyperref}

\ifr@index
  \RequirePackage{imakeidx}
  \makeindex
\fi

\makeatletter
%    \end{macrocode}
%
% \subsection{Internal helpers}
%
%    \begin{macrocode}
\newcommand{\r@trimid}[1]{\zap@space#1 \@empty}
\newcommand{\r@warn}[1]{\PackageWarning{registerstyle}{#1}}
\newcommand{\r@indexperson}[1]{\ifr@index\index{#1}\fi}

\newcommand{\r@unknowngeneration}{0}

\newcommand{\r@csvadd}[2]{%
  \ifx#1\@empty
    \edef#1{#2}%
  \else
    \edef#1{#1,#2}%
  \fi
}
%    \end{macrocode}
%
% \subsection{Person registry}
%
% \begin{macro}{\r@rememberperson}
% Add a person id to the global person list if not already present.
%    \begin{macrocode}
\newcommand{\r@personlist}{}
\newcommand{\r@rememberperson}[1]{%
  \edef\r@id{\r@trimid{#1}}%
  \ifx\r@id\@empty\else
    \ifcsdef{r@registered@\r@id}{}{%
      \csdef{r@registered@\r@id}{1}%
      \r@csvadd{\r@personlist}{\r@id}%
    }%
  \fi
}
%    \end{macrocode}
% \end{macro}
%
% \begin{macro}{\r@ensureperson}
% Conditional that branches on whether a person is defined.
%    \begin{macrocode}
\newcommand{\r@ensureperson}[1]{%
  \ifcsdef{r@name@#1}{%
    \@firstoftwo
  }{%
    \r@warn{Undefined person `#1'}%
    \@secondoftwo
  }%
}
%    \end{macrocode}
% \end{macro}
%
% \begin{macro}{\r@definepersoninternal}
% Internal 8-argument person definition used by |\defineperson|.
%    \begin{macrocode}
\newcommand{\r@definepersoninternal}[8]{%
  \edef\r@id{\r@trimid{#1}}%
  \ifx\r@id\@empty
    \r@warn{Empty person id ignored}%
  \else
    \r@rememberperson{\r@id}%
    \ifcsdef{r@name@\r@id}{\r@warn{Redefining person `\r@id'}}{}%
    \csgdef{r@name@\r@id}{#3}%
    \csgdef{r@parents@\r@id}{#2}%
    \csgdef{r@preferred@\r@id}{#4}%
    \csgdef{r@index@\r@id}{#5}%
    \csgdef{r@lineageparent@\r@id}{#6}%
    \csgdef{r@sex@\r@id}{#7}%
    \csgdef{r@birth@\r@id}{#8}%
    \csgdef{r@birthplace@\r@id}{}%
    \csgdef{r@death@\r@id}{}%
    \csgdef{r@deathplace@\r@id}{}%
    \csgdef{r@baptism@\r@id}{}%
    \csgdef{r@baptismplace@\r@id}{}%
    \csgdef{r@burial@\r@id}{}%
    \csgdef{r@burialplace@\r@id}{}%
    \csgdef{r@occupation@\r@id}{}%
    \csgdef{r@marriagecount@\r@id}{0}%
    \csundef{r@generation@\r@id}%
    \csundef{r@generation@working@\r@id}%
  \fi
}
%    \end{macrocode}
% \end{macro}
%
% \begin{macro}{\defineperson}
% Legacy positional person definition.
%    \begin{macrocode}
\NewDocumentCommand{\defineperson}{mmmm o}{%
  \IfNoValueTF{#5}{%
    \r@definepersoninternal{#1}{#2}{#3}{#4}{#3}{}{}{}%
  }{%
    \r@definepersoninternal{#1}{#2}{#3}{#4}{#5}{}{}{}%
  }%
}
%    \end{macrocode}
% \end{macro}
%
% \begin{macro}{\definepersonkv}
% Key-value person definition.
%    \begin{macrocode}
\define@key{r@person}{name}{\def\r@kv@name{#1}}
\define@key{r@person}{preferred}{\def\r@kv@preferred{#1}}
\define@key{r@person}{parents}{\def\r@kv@parents{#1}}
\define@key{r@person}{index}{\def\r@kv@index{#1}}
\define@key{r@person}{lineage-parent}{\def\r@kv@lineageparent{#1}}
\define@key{r@person}{sex}{\def\r@kv@sex{#1}}
\define@key{r@person}{birth}{\def\r@kv@birth{#1}}
\define@key{r@person}{birth-place}{\def\r@kv@birthplace{#1}}
\define@key{r@person}{death}{\def\r@kv@death{#1}}
\define@key{r@person}{death-place}{\def\r@kv@deathplace{#1}}
\define@key{r@person}{baptism}{\def\r@kv@baptism{#1}}
\define@key{r@person}{baptism-place}{\def\r@kv@baptismplace{#1}}
\define@key{r@person}{burial}{\def\r@kv@burial{#1}}
\define@key{r@person}{burial-place}{\def\r@kv@burialplace{#1}}
\define@key{r@person}{occupation}{\def\r@kv@occupation{#1}}
\define@key{r@person}{marriage-date}{\def\r@kv@marriagedate{#1}}
\define@key{r@person}{marriage-place}{\def\r@kv@marriageplace{#1}}

\NewDocumentCommand{\definepersonkv}{mm}{%
  \def\r@kv@name{}%
  \def\r@kv@preferred{}%
  \def\r@kv@parents{}%
  \def\r@kv@index{}%
  \def\r@kv@lineageparent{}%
  \def\r@kv@sex{}%
  \def\r@kv@birth{}%
  \def\r@kv@birthplace{}%
  \def\r@kv@death{}%
  \def\r@kv@deathplace{}%
  \def\r@kv@baptism{}%
  \def\r@kv@baptismplace{}%
  \def\r@kv@burial{}%
  \def\r@kv@burialplace{}%
  \def\r@kv@occupation{}%
  \def\r@kv@marriagedate{}%
  \def\r@kv@marriageplace{}%
  \setkeys{r@person}{#2}%
  \edef\r@id{\r@trimid{#1}}%
  \ifx\r@id\@empty
    \r@warn{Empty person id ignored}%
  \else
    \r@rememberperson{\r@id}%
    \ifx\r@kv@name\@empty
      \r@warn{Person `\r@id' defined without a name}%
    \fi
    \ifx\r@kv@index\@empty
      \let\r@kv@index\r@kv@name
    \fi
    \ifcsdef{r@name@\r@id}{\r@warn{Redefining person `\r@id'}}{}%
    \expandafter\global\expandafter\let\csname r@name@\r@id\endcsname\r@kv@name
    \expandafter\global\expandafter\let\csname r@parents@\r@id\endcsname\r@kv@parents
    \expandafter\global\expandafter\let\csname r@preferred@\r@id\endcsname\r@kv@preferred
    \expandafter\global\expandafter\let\csname r@index@\r@id\endcsname\r@kv@index
    \expandafter\global\expandafter\let\csname r@lineageparent@\r@id\endcsname\r@kv@lineageparent
    \expandafter\global\expandafter\let\csname r@sex@\r@id\endcsname\r@kv@sex
    \expandafter\global\expandafter\let\csname r@birth@\r@id\endcsname\r@kv@birth
    \expandafter\global\expandafter\let\csname r@birthplace@\r@id\endcsname\r@kv@birthplace
    \expandafter\global\expandafter\let\csname r@death@\r@id\endcsname\r@kv@death
    \expandafter\global\expandafter\let\csname r@deathplace@\r@id\endcsname\r@kv@deathplace
    \expandafter\global\expandafter\let\csname r@baptism@\r@id\endcsname\r@kv@baptism
    \expandafter\global\expandafter\let\csname r@baptismplace@\r@id\endcsname\r@kv@baptismplace
    \expandafter\global\expandafter\let\csname r@burial@\r@id\endcsname\r@kv@burial
    \expandafter\global\expandafter\let\csname r@burialplace@\r@id\endcsname\r@kv@burialplace
    \expandafter\global\expandafter\let\csname r@occupation@\r@id\endcsname\r@kv@occupation
    \csgdef{r@marriagecount@\r@id}{0}%
    \ifx\r@kv@marriagedate\@empty\else
      \csgdef{r@marriagecount@\r@id}{1}%
      \expandafter\global\expandafter\let\csname r@marriagedate@\r@id @1\endcsname\r@kv@marriagedate
      \csgdef{r@marriageplace@\r@id @1}{}%
    \fi
    \ifx\r@kv@marriageplace\@empty\else
      \ifcsdef{r@marriagedate@\r@id @1}{%
        \expandafter\global\expandafter\let\csname r@marriageplace@\r@id @1\endcsname\r@kv@marriageplace
      }{%
        \csgdef{r@marriagecount@\r@id}{1}%
        \csgdef{r@marriagedate@\r@id @1}{}%
        \expandafter\global\expandafter\let\csname r@marriageplace@\r@id @1\endcsname\r@kv@marriageplace
      }%
    \fi
    \csundef{r@generation@\r@id}%
    \csundef{r@generation@working@\r@id}%
  \fi
}
%    \end{macrocode}
% \end{macro}
%
% \subsection{Multiple marriage support}
%
% \begin{macro}{\addmarriage}
%    \begin{macrocode}
\define@key{r@marriage}{date}{\def\r@mkv@date{#1}}
\define@key{r@marriage}{place}{\def\r@mkv@place{#1}}
\define@key{r@marriage}{spouse}{\def\r@mkv@spouse{#1}}

\NewDocumentCommand{\addmarriage}{mm}{%
  \def\r@mkv@date{}%
  \def\r@mkv@place{}%
  \def\r@mkv@spouse{}%
  \setkeys{r@marriage}{#2}%
  \edef\r@id{\r@trimid{#1}}%
  \ifx\r@id\@empty
    \r@warn{Cannot add marriage to empty person id}%
  \else
    \ifcsdef{r@name@\r@id}{%
      \edef\r@mcnt{\csuse{r@marriagecount@\r@id}}%
      \edef\r@newmcnt{\the\numexpr\r@mcnt+1\relax}%
      \csgdef{r@marriagecount@\r@id}{\r@newmcnt}%
      \expandafter\global\expandafter\let\csname r@marriagedate@\r@id @\r@newmcnt\endcsname\r@mkv@date
      \expandafter\global\expandafter\let\csname r@marriageplace@\r@id @\r@newmcnt\endcsname\r@mkv@place
      \expandafter\global\expandafter\let\csname r@marriagespouse@\r@id @\r@newmcnt\endcsname\r@mkv@spouse
    }{%
      \r@warn{Cannot add marriage to undefined person `\r@id'}%
    }%
  \fi
}
%    \end{macrocode}
% \end{macro}
%
% \begin{macro}{\getmarriagecount}
% \begin{macro}{\getmarriagedate}
% \begin{macro}{\getmarriageplace}
% \begin{macro}{\getmarriagespouse}
%    \begin{macrocode}
\newcommand{\getmarriagecount}[1]{%
  \edef\r@id{\r@trimid{#1}}%
  \ifcsdef{r@marriagecount@\r@id}{\csuse{r@marriagecount@\r@id}}{0}%
}

\newcommand{\getmarriagedate}[2]{%
  \edef\r@id{\r@trimid{#1}}%
  \ifcsdef{r@marriagedate@\r@id @#2}{%
    \ifcsempty{r@marriagedate@\r@id @#2}{%
      \r@warn{No date stored for marriage #2 of `\r@id'}%
      \textbf{[marriage date unknown]}%
    }{%
      \edef\r@mdate{\csuse{r@marriagedate@\r@id @#2}}%
      \expandafter\registerdate\expandafter{\r@mdate}%
    }%
  }{%
    \r@warn{No marriage #2 for `\r@id'}%
    \textbf{[no marriage #2]}%
  }%
}

\newcommand{\getmarriageplace}[2]{%
  \edef\r@id{\r@trimid{#1}}%
  \ifcsdef{r@marriageplace@\r@id @#2}{%
    \ifcsempty{r@marriageplace@\r@id @#2}{%
      \r@warn{No place stored for marriage #2 of `\r@id'}%
      \textbf{[marriage place unknown]}%
    }{%
      \csuse{r@marriageplace@\r@id @#2}%
    }%
  }{%
    \r@warn{No marriage #2 for `\r@id'}%
    \textbf{[no marriage #2]}%
  }%
}

\newcommand{\getmarriagespouse}[2]{%
  \edef\r@id{\r@trimid{#1}}%
  \ifcsdef{r@marriagespouse@\r@id @#2}{%
    \ifcsempty{r@marriagespouse@\r@id @#2}{}{%
      \edef\r@spid{\csuse{r@marriagespouse@\r@id @#2}}%
      \getpersonname{\r@spid}%
    }%
  }{}%
}
%    \end{macrocode}
% \end{macro}
% \end{macro}
% \end{macro}
% \end{macro}
%
% \subsection{Name retrieval}
%
% \begin{macro}{\getpersonname}
% \begin{macro}{\getpersonnameNI}
% \begin{macro}{\getpersonpreferred}
% \begin{macro}{\getpersonabbrevNI}
% \begin{macro}{\shortname}
%    \begin{macrocode}
\newcommand{\getpersonname}[1]{%
  \edef\r@id{\r@trimid{#1}}%
  \ifcsdef{r@name@\r@id}{%
    \leavevmode
    \textsc{\csuse{r@name@\r@id}}%
    \r@indexperson{\csuse{r@index@\r@id}}%
  }{%
    \r@warn{Undefined person `\r@id'}%
    \textbf{[undefined \r@id]}%
  }%
}

\newcommand{\getpersonnameNI}[1]{%
  \edef\r@id{\r@trimid{#1}}%
  \ifcsdef{r@name@\r@id}{%
    \textsc{\csuse{r@name@\r@id}}%
  }{%
    \r@warn{Undefined person `\r@id'}%
    \textbf{[undefined \r@id]}%
  }%
}

\newcommand{\getpersonpreferred}[1]{%
  \edef\r@id{\r@trimid{#1}}%
  \ifcsdef{r@preferred@\r@id}{%
    \ifcsempty{r@preferred@\r@id}{}{\textsc{\csuse{r@preferred@\r@id}}}%
  }{}%
}

\newcommand{\getpersonabbrevNI}[1]{%
  \edef\r@id{\r@trimid{#1}}%
  \ifcsdef{r@preferred@\r@id}{%
    \ifcsempty{r@preferred@\r@id}{%
      \getpersonnameNI{\r@id}%
    }{%
      \textsc{\csuse{r@preferred@\r@id}}%
    }%
  }{%
    \getpersonnameNI{\r@id}%
  }%
}

\newcommand{\shortname}[1]{\getpersonpreferred{#1}}
%    \end{macrocode}
% \end{macro}
% \end{macro}
% \end{macro}
% \end{macro}
% \end{macro}
%
% \subsection{Vital record accessors}
%
% \begin{macro}{\r@getvitaldate}
% \begin{macro}{\r@getvitalplace}
% Shared helpers for date and place retrieval.
%    \begin{macrocode}
\newcommand{\r@getvitaldate}[3]{%
  \edef\r@id{\r@trimid{#1}}%
  \ifcsdef{r@#2@\r@id}{%
    \ifcsempty{r@#2@\r@id}{%
      \r@warn{No #3 date stored for `\r@id'}%
      \textbf{[#3 date unknown for \r@id]}%
    }{%
      \edef\r@thedate{\csuse{r@#2@\r@id}}%
      \expandafter\registerdate\expandafter{\r@thedate}%
    }%
  }{%
    \r@warn{Undefined person `\r@id' while retrieving #3 date}%
    \textbf{[undefined \r@id]}%
  }%
}

\newcommand{\r@getvitalplace}[3]{%
  \edef\r@id{\r@trimid{#1}}%
  \ifcsdef{r@#2@\r@id}{%
    \ifcsempty{r@#2@\r@id}{%
      \r@warn{No #3 place stored for `\r@id'}%
      \textbf{[#3 place unknown for \r@id]}%
    }{%
      \csuse{r@#2@\r@id}%
    }%
  }{%
    \r@warn{Undefined person `\r@id' while retrieving #3 place}%
    \textbf{[undefined \r@id]}%
  }%
}

\newcommand{\getpersonbirthdate}[1]{\r@getvitaldate{#1}{birth}{birth}}
\newcommand{\getpersonbirthplace}[1]{\r@getvitalplace{#1}{birthplace}{birth}}
\newcommand{\getpersondeathdate}[1]{\r@getvitaldate{#1}{death}{death}}
\newcommand{\getpersondeathplace}[1]{\r@getvitalplace{#1}{deathplace}{death}}
\newcommand{\getpersonbaptismdate}[1]{\r@getvitaldate{#1}{baptism}{baptism}}
\newcommand{\getpersonbaptismplace}[1]{\r@getvitalplace{#1}{baptismplace}{baptism}}
\newcommand{\getpersonburialdate}[1]{\r@getvitaldate{#1}{burial}{burial}}
\newcommand{\getpersonburialplace}[1]{\r@getvitalplace{#1}{burialplace}{burial}}
%    \end{macrocode}
% \end{macro}
% \end{macro}
%
% \begin{macro}{\getpersonoccupation}
% \begin{macro}{\getpersonsex}
%    \begin{macrocode}
\newcommand{\getpersonoccupation}[1]{%
  \edef\r@id{\r@trimid{#1}}%
  \ifcsdef{r@occupation@\r@id}{%
    \ifcsempty{r@occupation@\r@id}{%
      \r@warn{No occupation stored for `\r@id'}%
      \textbf{[occupation unknown for \r@id]}%
    }{%
      \csuse{r@occupation@\r@id}%
    }%
  }{%
    \r@warn{Undefined person `\r@id' while retrieving occupation}%
    \textbf{[undefined \r@id]}%
  }%
}

\newcommand{\getpersonsex}[1]{%
  \edef\r@id{\r@trimid{#1}}%
  \ifcsdef{r@sex@\r@id}{\csuse{r@sex@\r@id}}{}%
}
%    \end{macrocode}
% \end{macro}
% \end{macro}
%
% \subsection{Gender-aware pronouns}
%
% \begin{macro}{\r@pronoun}
% Internal dispatcher: selects masculine, feminine, or neutral form.
%    \begin{macrocode}
\newcommand{\r@pronoun}[4]{%
  \edef\r@id{\r@trimid{#1}}%
  \ifcsdef{r@sex@\r@id}{%
    \edef\r@thesex{\csuse{r@sex@\r@id}}%
    \def\r@male{M}\def\r@female{F}%
    \ifx\r@thesex\r@male #2%
    \else\ifx\r@thesex\r@female #3%
    \else #4%
    \fi\fi
  }{#4}%
}

\newcommand{\heshethey}[1]{\r@pronoun{#1}{he}{she}{they}}
\newcommand{\himherthem}[1]{\r@pronoun{#1}{him}{her}{them}}
\newcommand{\hishertheirargs}[1]{\r@pronoun{#1}{his}{her}{their}}
\newcommand{\HeSheThey}[1]{\r@pronoun{#1}{He}{She}{They}}
\newcommand{\HimHerThem}[1]{\r@pronoun{#1}{Him}{Her}{Them}}
\newcommand{\HisHerTheir}[1]{\r@pronoun{#1}{His}{Her}{Their}}
\newcommand{\himselfherselfthemself}[1]{\r@pronoun{#1}{himself}{herself}{themself}}
%    \end{macrocode}
% \end{macro}
%
% \subsection{Approximate dates}
%
%    \begin{macrocode}
\newcommand{\circa}[1]{c.\,#1}
\newcommand{\before}[1]{bef.\ #1}
\newcommand{\after}[1]{aft.\ #1}
\newcommand{\between}[2]{bet.\ #1 and #2}
\newcommand{\aboutdate}[1]{abt.\ #1}
\newcommand{\probably}[1]{prob.\ #1}
\newcommand{\say}[1]{say #1}
%    \end{macrocode}
%
% \subsection{Generation computation}
%
% \begin{macro}{\r@computegeneration}
% Recursive computation with cycle detection via a ``working'' flag.
%    \begin{macrocode}
\newcommand{\r@computegeneration}[1]{%
  \begingroup
    \edef\r@target{\r@trimid{#1}}%
    \ifcsdef{r@generation@\r@target}{%
    }{%
      \ifcsdef{r@generation@working@\r@target}{%
        \r@warn{Cycle detected while computing generation for `\r@target'}%
        \expandafter\xdef\csname r@generation@\r@target\endcsname{\r@unknowngeneration}%
      }{%
        \csdef{r@generation@working@\r@target}{1}%
        \ifcsdef{r@parents@\r@target}{%
          \edef\r@parents{\csuse{r@parents@\r@target}}%
          \ifx\r@parents\@empty
            \expandafter\xdef\csname r@generation@\r@target\endcsname{1}%
          \else
            \def\r@max{0}%
            \def\r@checkparent##1{%
              \edef\r@pid{\r@trimid{##1}}%
              \ifx\r@pid\@empty
                \r@warn{Empty parent id listed for `\r@target'}%
              \else
                \ifcsdef{r@parents@\r@pid}{%
                  \ifcsdef{r@generation@\r@pid}{}{\r@computegeneration{\r@pid}}%
                  \edef\r@pg{\csuse{r@generation@\r@pid}}%
                  \ifnum\r@pg>\r@max\relax \xdef\r@max{\r@pg}\fi
                }{%
                  \r@warn{Undefined parent `\r@pid' listed for `\r@target'}%
                  \expandafter\xdef\csname r@generation@\r@pid\endcsname{\r@unknowngeneration}%
                }%
              \fi
            }%
            \expandafter\forcsvlist\expandafter\r@checkparent\expandafter{\r@parents}%
            \edef\r@this{\the\numexpr\r@max+1\relax}%
            \expandafter\xdef\csname r@generation@\r@target\endcsname{\r@this}%
          \fi
        }{%
          \r@warn{Undefined person `\r@target'; assigning generation \r@unknowngeneration}%
          \expandafter\xdef\csname r@generation@\r@target\endcsname{\r@unknowngeneration}%
        }%
        \csundef{r@generation@working@\r@target}%
      }%
    }%
  \endgroup
}
%    \end{macrocode}
% \end{macro}
%
% \begin{macro}{\getgeneration}
%    \begin{macrocode}
\newcommand{\getgeneration}[1]{%
  \begingroup
    \edef\r@target{\r@trimid{#1}}%
    \r@computegeneration{\r@target}%
    \xdef\r@generationresult{\csuse{r@generation@\r@target}}%
  \endgroup
  \r@generationresult
}
%    \end{macrocode}
% \end{macro}
%
% \subsection{Ancestor chain formatting}
%
% \begin{macro}{\r@collectancestors}
% \begin{macro}{\r@addoneancestor}
%    \begin{macrocode}
\newcommand{\r@alist}{}
\newif\ifr@ancestorchain@lineage

\newcommand{\r@collectancestors}[1]{%
  \edef\r@id{\r@trimid{#1}}%
  \ifcsdef{r@ancestor@seen@\r@id}{%
    \r@warn{Cycle detected while collecting ancestors for `\r@id'}%
  }{%
    \csdef{r@ancestor@seen@\r@id}{1}%
    \ifcsdef{r@parents@\r@id}{%
      \ifcsdef{r@lineageparent@\r@id}{%
        \ifcsempty{r@lineageparent@\r@id}{%
          \edef\r@plist{\csuse{r@parents@\r@id}}%
          \ifx\r@plist\@empty\else
            \expandafter\forcsvlist\expandafter\r@addoneancestor\expandafter{\r@plist}%
          \fi
        }{%
          \r@ancestorchain@lineagetrue
          \edef\r@lineagepid{\csuse{r@lineageparent@\r@id}}%
          \expandafter\r@addoneancestor\expandafter{\r@lineagepid}%
        }%
      }{%
        \edef\r@plist{\csuse{r@parents@\r@id}}%
        \ifx\r@plist\@empty\else
          \expandafter\forcsvlist\expandafter\r@addoneancestor\expandafter{\r@plist}%
        \fi
      }%
    }{%
      \r@warn{Undefined person `\r@id' while collecting ancestors}%
    }%
  }%
}
\newcommand{\r@addoneancestor}[1]{%
  \edef\r@pid{\r@trimid{#1}}%
  \ifx\r@pid\@empty
    \r@warn{Empty parent id in ancestor chain}%
  \else
    \ifcsdef{r@parents@\r@pid}{%
      \r@csvadd{\r@alist}{\r@pid}%
      \r@collectancestors{\r@pid}%
    }{%
      \r@warn{Undefined parent `\r@pid' in ancestor chain}%
      \r@csvadd{\r@alist}{\r@pid}%
    }%
  \fi
}
%    \end{macrocode}
% \end{macro}
% \end{macro}
%
% \begin{macro}{\formatparents}
%    \begin{macrocode}
\newif\ifr@abbrevparents
\ifr@option@abbrevparents
  \r@abbrevparentstrue
\else
  \r@abbrevparentsfalse
\fi

\newcount\r@pcnt
\newcount\r@i

\newcommand{\r@namegensNI}[1]{%
  \mbox{\getpersonnameNI{#1}\textsuperscript{\getgeneration{#1}}}%
}
\newcommand{\r@abbrevnamegensNI}[1]{%
  \mbox{\getpersonabbrevNI{#1}\textsuperscript{\getgeneration{#1}}}%
}

\newcommand{\formatparents}[1]{%
  \begingroup
    \def\r@alist{}%
    \r@ancestorchain@lineagefalse
    \r@collectancestors{#1}%
    \edef\r@plist{\r@alist}%
    \ifx\r@plist\@empty
    \else
      \r@pcnt=0\relax
      \def\r@countone##1{\advance\r@pcnt by 1\relax}%
      \expandafter\forcsvlist\expandafter\r@countone\expandafter{\r@plist}%
      \unskip\space(%
        \r@i=0\relax
        \def\r@printone##1{%
          \edef\r@pid{\r@trimid{##1}}%
          \advance\r@i by 1\relax
          \ifnum\r@i>1
            \ifr@ancestorchain@lineage
              ,\space
            \else
              \ifnum\r@i=\r@pcnt
                \space and\space
              \else
                ,\space
              \fi
            \fi
          \fi
          \ifr@abbrevparents
            \r@abbrevnamegensNI{\r@pid}%
          \else
            \r@namegensNI{\r@pid}%
          \fi
        }%
        \expandafter\forcsvlist\expandafter\r@printone\expandafter{\r@plist}%
      )%
    \fi
  \endgroup
}
%    \end{macrocode}
% \end{macro}
%
% \subsection{Main register entries}
%
% \begin{macro}{\mainperson}
%    \begin{macrocode}
\newcommand{\registerpersonsep}{1em}
\newif\ifr@afterheading
\r@afterheadingfalse

\pretocmd{\section}{\global\r@afterheadingtrue}{}{}
\pretocmd{\subsection}{\global\r@afterheadingtrue}{}{}
\pretocmd{\subsubsection}{\global\r@afterheadingtrue}{}{}
\pretocmd{\paragraph}{\global\r@afterheadingtrue}{}{}
\pretocmd{\subparagraph}{\global\r@afterheadingtrue}{}{}

\newcounter{registerperson}

\NewDocumentCommand{\mainperson}{m}{%
  \edef\r@id{\r@trimid{#1}}%
  \par
  \ifr@afterheading
    \global\r@afterheadingfalse
  \else
    \addvspace{\registerpersonsep}%
  \fi
  \refstepcounter{registerperson}%
  \csgdef{r@mainentry@\r@id}{\theregisterperson}%
  \label{person:\r@id}%
  \noindent
  \textbf{\theregisterperson.\ }%
  \getpersonname{\r@id}\textsuperscript{\getgeneration{\r@id}}%
  \formatparents{\r@id}\xspace
}
%    \end{macrocode}
% \end{macro}
%
% NGSQ numbering override.
%    \begin{macrocode}
\ifr@option@ngsq
  \renewcommand{\theregisterperson}{\arabic{registerperson}}
\fi
%    \end{macrocode}
%
% \subsection{Cross-references}
%
% \begin{macro}{\xrefperson}
% \begin{macro}{\childname}
%    \begin{macrocode}
\NewDocumentCommand{\xrefperson}{m}{%
  \edef\r@id{\r@trimid{#1}}%
  \hyperref[person:\r@id]{%
    \getpersonname{\r@id}\textsuperscript{\getgeneration{\r@id}}%
    \space(\#\ref{person:\r@id})%
  }%
}

\NewDocumentCommand{\childname}{m}{%
  \edef\r@id{\r@trimid{#1}}%
  \getpersonnameNI{\r@id}\textsuperscript{\getgeneration{\r@id}}%
}
%    \end{macrocode}
% \end{macro}
% \end{macro}
%
% \subsection{Child lists}
%
% \begin{environment}{children}
% \begin{environment}{childrenof}
%    \begin{macrocode}
\newif\ifr@childrenof
\r@childrenoffalse
\newcommand{\r@childrenof@a}{}
\newcommand{\r@childrenof@b}{}

\NewDocumentEnvironment{children}{ O{} }{%
  \begin{enumerate}[label=\roman*., leftmargin=2em, itemsep=0pt, topsep=0pt,#1]%
}{\end{enumerate}}

\NewDocumentEnvironment{childrenof}{ m m O{} }{%
  \begingroup
  \edef\r@childrenof@a{\r@trimid{#1}}%
  \edef\r@childrenof@b{\r@trimid{#2}}%
  \r@childrenoftrue
  \begin{children}[#3]%
}{%
  \end{children}%
  \endgroup
}
%    \end{macrocode}
% \end{environment}
% \end{environment}
%
% \begin{macro}{\r@checkchildparents}
% \begin{macro}{\child}
% \begin{macro}{\childref}
%    \begin{macrocode}
\newcommand{\r@checkchildparents}[1]{%
  \ifr@childrenof
    \edef\r@childid{\r@trimid{#1}}%
    \ifcsdef{r@parents@\r@childid}{%
      \edef\r@childparents{\csuse{r@parents@\r@childid}}%
      \edef\r@expectedparents{\r@childrenof@a,\r@childrenof@b}%
      \edef\r@expectedparentsrev{\r@childrenof@b,\r@childrenof@a}%
      \ifx\r@childparents\r@expectedparents
      \else
        \ifx\r@childparents\r@expectedparentsrev
        \else
          \r@warn{Child `\r@childid' listed under parents `\r@expectedparents' but has parents `\r@childparents'}%
        \fi
      \fi
    }{%
      \r@warn{Undefined child `\r@childid' listed in childrenof}%
    }%
  \fi
}

\NewDocumentCommand{\child}{m m}{%
  \r@checkchildparents{#1}%
  \item \childname{#1}%
  \IfBlankF{#2}{, #2}%
}

\NewDocumentCommand{\childref}{m m}{%
  \r@checkchildparents{#1}%
  \item \textbf{+}\space\childname{#1}%
  \IfBlankF{#2}{, #2}%
}
%    \end{macrocode}
% \end{macro}
% \end{macro}
% \end{macro}
%
% \subsection{Source and note support}
%
% \begin{macro}{\registersource}
% \begin{macro}{\registerendnote}
%    \begin{macrocode}
\NewDocumentCommand{\registersource}{m}{%
  \footnote{#1}%
}

\NewDocumentCommand{\registerendnote}{m}{%
  \@ifpackageloaded{endnotes}{\endnote{#1}}{\footnote{#1}}%
}
%    \end{macrocode}
% \end{macro}
% \end{macro}
%
% \subsection{Automatic generation headings}
%
% \begin{macro}{\generationheading}
%    \begin{macrocode}
\newcommand{\r@generationword}[1]{%
  \ifcase#1 Zeroth\or First\or Second\or Third\or Fourth\or Fifth\or
    Sixth\or Seventh\or Eighth\or Ninth\or Tenth\or Eleventh\or
    Twelfth\or Thirteenth\or Fourteenth\or Fifteenth\or Sixteenth\or
    Seventeenth\or Eighteenth\or Nineteenth\or Twentieth\else
    #1th\fi
}

\newcommand{\generationheading}[1]{%
  \edef\r@id{\r@trimid{#1}}%
  \r@computegeneration{\r@id}%
  \edef\r@gen{\csuse{r@generation@\r@id}}%
  \par\medskip\noindent
  \textbf{\r@generationword{\r@gen}\space Generation}%
  \par\nobreak
}
%    \end{macrocode}
% \end{macro}
%
% \subsection{Date formatting}
%
% \begin{macro}{\registerdate}
% \begin{macro}{\registerdatestyle}
%    \begin{macrocode}
\newcommand{\registerdatestyle}[1]{%
  \DTMsetdatestyle{#1}%
}

\AtBeginDocument{%
  \ifdefempty{\r@datestyle}{}{\DTMsetdatestyle{\r@datestyle}}%
}

\ExplSyntaxOn
\cs_new_protected:Npn \registerstyle_date:n #1
  {
    \regex_match:nnTF { \A \d{4} - \d{2} - \d{2} \Z } {#1}
      { \DTMdate{#1} }
      {
        \PackageWarning{registerstyle}{Date~`#1'~is~not~in~YYYY-MM-DD~form}
        #1
      }
  }

\NewDocumentCommand{\registerdate}{m}{\registerstyle_date:n{#1}}
\ExplSyntaxOff
%    \end{macrocode}
% \end{macro}
% \end{macro}
%
% \subsection{Date validation}
%
%    \begin{macrocode}
\ExplSyntaxOn
\cs_generate_variant:Nn \regex_match:nnF { neF }
\cs_new_protected:Npn \registerstyle_validate_iso_date:nn #1 #2
  {
    \tl_if_blank:nF {#2}
      {
        \regex_match:neF { \A \d{4} - \d{2} - \d{2} \Z } {#2}
          { \PackageWarning{registerstyle}{#1~date~`#2'~is~not~in~YYYY-MM-DD~form} }
      }
  }
\NewDocumentCommand{\r@validateisodate}{m m}
  {\registerstyle_validate_iso_date:nn {#1} {#2}}
\ExplSyntaxOff
%    \end{macrocode}
%
% \subsection{Validation commands}
%
% \begin{macro}{\checkperson}
% \begin{macro}{\checkregister}
%    \begin{macrocode}
\newcommand{\r@checkdatefield}[3]{%
  \ifcsdef{r@#3@#1}{%
    \edef\r@datevalue{\csuse{r@#3@#1}}%
    \ifx\r@datevalue\@empty\else
      \r@validateisodate{#2 for `#1'}{\r@datevalue}%
    \fi
  }{}%
}

\newcommand{\r@checkplacewithoutdate}[5]{%
  \ifcsdef{r@#2@#1}{%
    \ifcsempty{r@#2@#1}{}{%
      \ifcsdef{r@#3@#1}{%
        \ifcsempty{r@#3@#1}{\r@warn{Person `#1' has #4 but no #5}}{}%
      }{\r@warn{Person `#1' has #4 but no #5}}%
    }%
  }{}%
}

\newcommand{\r@checkparentlist}[1]{%
  \ifcsdef{r@parents@#1}{%
    \edef\r@plist{\csuse{r@parents@#1}}%
    \ifx\r@plist\@empty\else
      \def\r@checkoneparent##1{%
        \edef\r@pid{\r@trimid{##1}}%
        \ifx\r@pid\@empty
          \r@warn{Person `#1' has an empty parent id}%
        \else
          \ifcsdef{r@name@\r@pid}{}{\r@warn{Person `#1' lists undefined parent `\r@pid'}}%
        \fi
      }%
      \expandafter\forcsvlist\expandafter\r@checkoneparent\expandafter{\r@plist}%
    \fi
  }{\r@warn{Person `#1' has no parent field}}%
}

\NewDocumentCommand{\checkperson}{m}{%
  \begingroup
    \edef\r@id{\r@trimid{#1}}%
    \ifx\r@id\@empty
      \r@warn{Cannot check an empty person id}%
    \else
      \ifcsdef{r@name@\r@id}{%
        \ifcsempty{r@name@\r@id}{\r@warn{Person `\r@id' has an empty name}}{}%
        \r@checkparentlist{\r@id}%
        \r@checkdatefield{\r@id}{Birth}{birth}%
        \r@checkdatefield{\r@id}{Death}{death}%
        \r@checkdatefield{\r@id}{Baptism}{baptism}%
        \r@checkdatefield{\r@id}{Burial}{burial}%
        \r@checkplacewithoutdate{\r@id}{birthplace}{birth}{birth place}{birth date}%
        \r@checkplacewithoutdate{\r@id}{deathplace}{death}{death place}{death date}%
        \r@checkplacewithoutdate{\r@id}{baptismplace}{baptism}{baptism place}{baptism date}%
        \r@checkplacewithoutdate{\r@id}{burialplace}{burial}{burial place}{burial date}%
      }{%
        \r@warn{Cannot check undefined person `\r@id'}%
      }%
    \fi
  \endgroup
}

\NewDocumentCommand{\checkregister}{}{%
  \begingroup
    \ifx\r@personlist\@empty
      \r@warn{No registered persons to check}%
    \else
      \def\r@checkoneperson##1{\checkperson{##1}}%
      \expandafter\forcsvlist\expandafter\r@checkoneperson\expandafter{\r@personlist}%
    \fi
  \endgroup
}
%    \end{macrocode}
% \end{macro}
% \end{macro}
%
% \subsection{Shorthands}
%
%    \begin{macrocode}
\newcommand{\born}[1]{b.\ #1}
\newcommand{\died}[1]{d.\ #1}
\newcommand{\married}[1]{m.\ #1}
\newcommand{\baptized}[1]{bp.\ #1}
\newcommand{\buried}[1]{bur.\ #1}
\newcommand{\place}[1]{\textit{#1}}
\newcommand{\dated}[1]{#1}
%    \end{macrocode}
%
% \subsection{Register section environment}
%
% \begin{environment}{registersection}
%    \begin{macrocode}
\NewDocumentEnvironment{registersection}{ O{} }{%
  \begingroup
  \setlength{\parindent}{1.5em}
  \setlength{\parskip}{0pt}
  \setlist{nosep}
  \def\generation##1{\par\medskip\noindent\textbf{##1 Generation}\par\nobreak}%
}{%
  \par\endgroup
}
%    \end{macrocode}
% \end{environment}
%
% \subsection{Backward compatibility}
%
% Old public names retained as aliases for one release cycle.
%    \begin{macrocode}
\let\collectancestors\r@collectancestors
\let\addoneancestor\r@addoneancestor
\newif\ifabbrevparents
\ifr@abbrevparents\abbrevparentstrue\else\abbrevparentsfalse\fi

\makeatother
%</package>
%    \end{macrocode}
%
% \Finale
